\chapter{2008年广东卷（文）}

\section{单选题}

\begin{problem}
第二十九届夏季奥林匹克运动会将于 \(2008\) 年 \(8\) 月 \(8\) 日在北京举行. 若集合 \(A = \{\) 参加北京奥运会比赛的运动员 \(\}\)，集合 \(B = \{\) 参加北京奥运会比赛的男运动员 \(\}\)，集合 \(C = \{\) 参加北京奥运会比赛的女运动员 \(\}\)，则下列关系正确的是（\quad）

\choices
    {\(A \subseteq B\)}
    {\(B \subseteq C\)}
    {\(A \cap B = C\)}
    {\(B \cup C = A\)}
\end{problem}

\begin{answer}
D
\end{answer}

\begin{solution}
参加比赛的运动员不是男运动员就是女运动员，因此集合 \(A\) 由集合 \(B\) 与集合 \(C\) 的所有元素组成，即 \(B\cup C=A\)，故选 D.
\end{solution}

\begin{problem}
已知 \(0 < a < 2\)，复数 \(z = a + \i\)（\(\i\) 是虚数单位），则 \(|z|\) 的取值范围是（\quad）

\choices
    {\((1, \sqrt{3})\)}
    {\((1, \sqrt{5})\)}
    {\((1, 3)\)}
    {\((1, 5)\)}
\end{problem}

\begin{answer}
B
\end{answer}

\begin{solution}
由复数模的定义，\(|z|=\sqrt{a^2+1}\). 因为 \(0<a<2\)，所以 \(1<a^2+1<5\)，而平方根函数在正数集上单调递增，故 \(1<|z|<\sqrt{5}\)，选 B.
\end{solution}

\begin{problem}
已知平面向量 \(\bs{a} = (1, 2)\)，\(\bs{b} = (-2, m)\)，且 \(\bs{a} \parallel \bs{b}\)，则 \(2\bs{a} + 3\bs{b} =\)（\quad）

\choices
    {\((-2, -4)\)}
    {\((-3, -6)\)}
    {\((-4, -8)\)}
    {\((-5, -10)\)}
\end{problem}

\begin{answer}
C
\end{answer}

\begin{solution}
由 \(\bs{a}\parallel\bs{b}\)，存在实数 \(k\) 使 \(\bs{b}=k\bs{a}\). 比较第一坐标得 \(k=-2\)，从而 \(m=-4\)，即 \(\bs{b}=(-2,-4)\). 因此
\(2\bs{a}+3\bs{b}=2(1,2)+3(-2,-4)=(-4,-8)\)，故选 C.
\end{solution}

\begin{problem}
已知等差数列 \(\{a_n\}\) 的前 \(n\) 项和为 \(S_n\)，若 \(S_2 = 4\)，\(S_4 = 20\)，则该数列的公差 \(d =\)（\quad）

\choices
    {\(7\)}
    {\(6\)}
    {\(3\)}
    {\(2\)}
\end{problem}

\begin{answer}
C
\end{answer}

\begin{solution}
设首项为 \(a_1\)，公差为 \(d\). 由等差数列前几项的和，得
\(S_2=2a_1+d=4\)，\(S_4=4a_1+6d=20\). 用第二式减去第一式的两倍，得 \(4d=12\)，所以 \(d=3\)，故选 C.
\end{solution}

\begin{problem}
已知函数 \(f(x) = (1 + \cos 2x)\sin^2 x\)，\(x \in \R\)，则 \(f(x)\) 是（\quad）

\choices
    {最小正周期为 \(\pi\) 的奇函数}
    {最小正周期为 \(\pi\) 的偶函数}
    {最小正周期为 \(\frac{\pi}{2}\) 的奇函数}
    {最小正周期为 \(\frac{\pi}{2}\) 的偶函数}
\end{problem}

\begin{answer}
D
\end{answer}

\begin{solution}
利用 \(1+\cos2x=2\cos^2x\)，有
\(f(x)=2\sin^2x\cos^2x=\frac{1}{2}\sin^22x=\frac14(1-\cos4x)\). 由最后一个式子可知 \(f(-x)=f(x)\)，所以它是偶函数；又因为 \(\cos4x\) 的最小正周期为 \(\frac{\pi}{2}\)，且系数非零，\(f(x)\) 的最小正周期也是 \(\frac{\pi}{2}\)，故选 D.
\end{solution}

\begin{problem}
经过圆 \(x^2 + 2x + y^2 = 0\) 的圆心 \(C\) 且与直线 \(x + y = 0\) 垂直的直线方程是（\quad）

\choices
    {\(x - y + 1 = 0\)}
    {\(x - y - 1 = 0\)}
    {\(x + y - 1 = 0\)}
    {\(x + y + 1 = 0\)}
\end{problem}

\begin{answer}
A
\end{answer}

\begin{solution}
将圆的方程配方，得 \((x+1)^2+y^2=1\)，所以圆心为 \(C(-1,0)\). 直线 \(x+y=0\) 的斜率为 \(-1\)，与它垂直的直线斜率为 \(1\). 过点 \(C\) 的直线为 \(y=x+1\)，即 \(x-y+1=0\)，故选 A.
\end{solution}

\begin{problem}
将正三棱柱截去三个角（如图1所示 \(A\)，\(B\)，\(C\) 分别是 \(\triangle GHI\) 三边的中点）得到几何体如图2，则该几何体按图2所示方向的侧视图（或称左视图）为（\quad）

\bitmapfigure[max width=0.72\linewidth]{img/2008/guangdong_liberal/q7_fig1.png}

\choices
    {\choicebitmap[width=0.15\paperwidth]{img/2008/guangdong_liberal/q7_choice_a.png}}
    {\choicebitmap[width=0.15\paperwidth]{img/2008/guangdong_liberal/q7_choice_b.png}}
    {\choicebitmap[width=0.15\paperwidth]{img/2008/guangdong_liberal/q7_choice_c.png}}
    {\choicebitmap[width=0.15\paperwidth]{img/2008/guangdong_liberal/q7_choice_d.png}}
\end{problem}

\begin{answer}
A
\end{answer}

\begin{solution}
沿图2箭头方向作平行投影，重合的顶点只保留一个投影点. 根据图2中外轮廓棱的投影关系，侧视图的外轮廓是左侧竖直、上边水平、右侧倾斜、下边水平的四边形，且棱 \(EB\) 的投影成为其中的内部线段. 该图形与选项 A 相同，故选 A.
\end{solution}

\begin{problem}
命题“若函数 \(f(x) = \log_a x\)（\(a > 0\)，\(a \neq 1\)）在其定义域内是减函数，则 \(\log_a 2 < 0\)”的逆否命题是（\quad）

\choices
    {若 \(\log_a 2 < 0\)，则函数 \(f(x) = \log_a x\)（\(a > 0\)，\(a \neq 1\)）在其定义域内不是减函数}
    {若 \(\log_a 2 \geqslant 0\)，则函数 \(f(x) = \log_a x\)（\(a > 0\)，\(a \neq 1\)）在其定义域内不是减函数}
    {若 \(\log_a 2 < 0\)，则函数 \(f(x) = \log_a x\)（\(a > 0\)，\(a \neq 1\)）在其定义域内是减函数}
    {若 \(\log_a 2 \geqslant 0\)，则函数 \(f(x) = \log_a x\)（\(a > 0\)，\(a \neq 1\)）在其定义域内是减函数}
\end{problem}

\begin{answer}
B
\end{answer}

\begin{solution}
原命题是“若 \(p\)，则 \(q\)”，其中 \(p\) 表示函数在定义域内是减函数，\(q\) 表示 \(\log_a2<0\). 逆否命题为“若 \(\lnot q\)，则 \(\lnot p\)”，即若 \(\log_a2\geqslant0\)，则函数在定义域内不是减函数，故选 B.
\end{solution}

\begin{problem}
设 \(a \in \R\)，若函数 \(y = \e^x + ax\)（\(x \in \R\)）有大于零的极值点，则（\quad）

\choices
    {\(a < -1\)}
    {\(a > -1\)}
    {\(a > -\frac{1}{\e}\)}
    {\(a < -\frac{1}{\e}\)}
\end{problem}

\begin{answer}
A
\end{answer}

\begin{solution}
函数的导数为 \(y'=\e^x+a\). 若有极值点，则存在 \(x_0\) 使 \(\e^{x_0}+a=0\)，所以 \(a<0\)，且 \(x_0=\ln(-a)\). 极值点横坐标大于零等价于 \(\ln(-a)>0\)，即 \(-a>1\)，所以 \(a<-1\)，故选 A.
\end{solution}

\begin{problem}
设 \(a\)，\(b \in \R\)，若 \(a - |b| > 0\)，则下列不等式中正确的是（\quad）

\choices
    {\(b - a > 0\)}
    {\(a^3 + b^3 < 0\)}
    {\(b + a > 0\)}
    {\(a^2 - b^2 < 0\)}
\end{problem}

\begin{answer}
C
\end{answer}

\begin{solution}
由 \(a-|b|>0\) 得 \(a>|b|\geqslant0\)，所以 \(a>0\) 且 \(-a<b<a\). 因而 \(a+b>0\)，选项 C 正确；同时 \(b-a<0\)，故 A 错误；又
\(a^3+b^3=(a+b)(a^2-ab+b^2)>0\)，故 B 错误；最后 \(a^2-b^2=(a-|b|)(a+|b|)>0\)，故 D 错误，故选 C.
\end{solution}

\section{填空题}

\begin{problem}
为了调查某厂工人生产某种产品的能力，随机抽查了 \(20\) 位工人某天生产该产品的数量. 产品数量的分组区间为 \([45, 55)\)，\([55, 65)\)，\([65, 75)\)，\([75, 85)\)，\([85, 95)\). 由此得到频率分布直方图如图，则这 \(20\) 名工人中一天生产该产品数量在 \([55, 75)\) 的人数是\fillinblank{}.

\texfigure[max width=0.46\linewidth]{img_repaint/2008/guangdong_liberal/q11_fig1.tex}
\end{problem}

\begin{answer}
13
\end{answer}

\begin{solution}
由频率分布直方图可知，区间 \([55,65)\) 和 \([65,75)\) 的频率分别为 \(0.040\) 和 \(0.025\). 因而产品数量在 \([55,75)\) 内的频率为 \((0.040+0.025)\times10=0.65\)，人数为 \(20\times0.65=13\).
\end{solution}

\begin{problem}
若变量 \(x\)，\(y\) 满足 \(\begin{cases}
2x + y \leqslant 40,\\
x + 2y \leqslant 50,\\
x \geqslant 0,\\
y \geqslant 0.
\end{cases}\)，则 \(z = 3x + 2y\) 的最大值是\fillinblank{}.
\end{problem}

\begin{answer}
70
\end{answer}

\begin{solution}
可行域的顶点为 \((0,0)\)，\((20,0)\)，\((10,20)\)，\((0,25)\). 在这些顶点处，\(z=3x+2y\) 的值分别为 \(0\)，\(60\)，\(70\)，\(50\). 线性目标函数的最大值在可行域顶点处取得，所以最大值为 \(70\).
\end{solution}

\begin{problem}
阅读如图的程序框图. 若输入 \(m = 4\)，\(n = 3\)，则输出 \(a =\)\fillinblank{}，\(i =\)\fillinblank{}.

\texfigure[max width=0.52\linewidth]{img_repaint/2008/guangdong_liberal/q13_fig1.tex}
\end{problem}

\begin{answer}
12，3
\end{answer}

\begin{solution}
初始时 \(i=1\)，计算 \(a=mi=4\)，因 \(3\) 不能整除 \(4\)，故令 \(i=2\). 此时 \(a=mi=8\)，因 \(3\) 不能整除 \(8\)，故令 \(i=3\). 此时 \(a=mi=12\)，且 \(3\mid12\)，满足判断条件，程序输出 \(a=12\)，\(i=3\).
\end{solution}

\section{选考题：填空题}

\begin{problem}
已知曲线 \(C_1\)，\(C_2\) 的极坐标方程分别为 \(\rho \cos \theta = 3\)，\(\rho = 4 \cos \theta\)（\(\rho \geqslant 0\)，\(0 \leqslant \theta < \frac{\pi}{2}\)），则曲线 \(C_1\) 与 \(C_2\) 交点的极坐标为\fillinblank{}.
\end{problem}

\begin{answer}
\(\left(2\sqrt{3},\frac{\pi}{6}\right)\)
\end{answer}

\begin{solution}
由 \(\rho\cos\theta=3\) 及 \(\rho=4\cos\theta\)，得 \(4\cos^2\theta=3\). 由于 \(0\leqslant\theta<\frac{\pi}{2}\)，所以 \(\cos\theta=\frac{\sqrt3}{2}\)，从而 \(\theta=\frac{\pi}{6}\)，\(\rho=4\cos\theta=2\sqrt3\). 因此交点的极坐标为 \(\left(2\sqrt3,\frac\pi6\right)\).
\end{solution}

\begin{problem}
已知 \(PA\) 是圆 \(O\) 的切线，切点为 \(A\)，\(PA = 2\). \(AC\) 是圆 \(O\) 的直径，\(PC\) 与圆 \(O\) 交于点 \(B\)，\(PB = 1\)，则圆 \(O\) 的半径 \(R =\)\fillinblank{}.
\end{problem}

\begin{answer}
\(\sqrt{3}\)
\end{answer}

\begin{solution}
由切割线定理，\(PA^2=PB\cdot PC\)，所以 \(PC=\frac{PA^2}{PB}=4\). 因为 \(PA\) 是切线，故 \(PA\perp OA\)；又 \(A\)，\(O\)，\(C\) 共线，所以 \(PA\perp AC\). 在直角三角形 \(PAC\) 中，\(AC=2R\)，故
\(PC^2=PA^2+AC^2\)，即 \(16=4+4R^2\)，解得 \(R=\sqrt3\).
\end{solution}

\section{解答题}

\begin{problem}
已知函数 \(f(x) = A \sin (x + \varphi)\)（\(A > 0\)，\(0 < \varphi < \pi\)），\(x \in \R\) 的最大值是 1，其图象经过点 \(M\left( \frac{\pi}{3}, \frac{1}{2} \right)\).

\begin{enumerate}
\item 求 \(f(x)\) 的解析式；
\item 已知 \(\alpha\)，\(\beta \in \left( 0, \frac{\pi}{2} \right)\)，且 \(f(\alpha) = \frac{3}{5}\)，\(f(\beta) = \frac{12}{13}\)，求 \(f(\alpha - \beta)\) 的值.
\end{enumerate}
\end{problem}

\begin{solution}
\begin{enumerate}
\item 因为 \(A>0\)，函数 \(A\sin(x+\varphi)\) 的最大值为 \(A\)，所以 \(A=1\). 又因图象经过点 \(M\)，有 \(\sin\left(\frac\pi3+\varphi\right)=\frac12\). 由 \(0<\varphi<\pi\)，得 \(\frac\pi3<\frac\pi3+\varphi<\frac{4\pi}3\)，故 \(\frac\pi3+\varphi=\frac{5\pi}6\)，从而 \(\varphi=\frac\pi2\). 因此 \(f(x)=\sin\left(x+\frac\pi2\right)=\cos x\).

\item 由 \(f(x)=\cos x\) 及 \(\alpha\)，\(\beta\in\left(0,\frac\pi2\right)\)，得
\(\cos\alpha=\frac35\)，\(\cos\beta=\frac{12}{13}\)，以及 \(\sin\alpha=\frac45\)，\(\sin\beta=\frac5{13}\). 所以
\(f(\alpha-\beta)=\cos(\alpha-\beta)=\cos\alpha\cos\beta+\sin\alpha\sin\beta=\frac35\cdot\frac{12}{13}+\frac45\cdot\frac5{13}=\frac{56}{65}\).
\end{enumerate}
\end{solution}

\begin{problem}
某单位用 \(2160\) 万元购得一块空地. 计划在该地块上建造一栋至少 \(10\) 层，每层 \(2000\text{ 平方米}\) 的楼房. 经测算，如果将楼房建为 \(x\) 层（\(x \geqslant 10\)），则每平方米的平均建筑费用为 \(560 + 48x\)（单位：元）. 为了使楼房每平方米的平均综合费用最少，该楼房应建为多少层？

（注：平均综合费用 = 平均建筑费用 + 平均购地费用 = \(\frac{\text{购地总费用}}{\text{建筑总面积}}\)）
\end{problem}

\begin{solution}
设每平方米的平均综合费用为 \(y\) 元. 购地总费用为 \(2160\times10^4\) 元，建成 \(x\) 层时建筑总面积为 \(2000x\) 平方米，所以
\(y=560+48x+\frac{2160\times10^4}{2000x}=560+48x+\frac{10800}{x}\)，其中 \(x\geqslant10\). 由基本不等式，
\(48x+\frac{10800}{x}\geqslant2\sqrt{48x\cdot\frac{10800}{x}}=1440\)，等号成立当且仅当 \(48x=\frac{10800}{x}\)，即 \(x^2=225\). 由于 \(x\geqslant10\)，得 \(x=15\). 因此该楼房应建为 \(15\) 层.
\end{solution}

\begin{problem}
如图所示，四棱锥 \(P - ABCD\) 的底面 \(ABCD\) 是半径为 \(R\) 的圆的内接四边形，其中 \(BD\) 是圆的直径，\(\angle ABD = 60^\circ\)，\(\angle BDC = 45^\circ\). \(\triangle ADP \sim \triangle BAD\).

\begin{enumerate}
\item 求线段 \(PD\) 的长；
\item 若 \(PC = \sqrt{11}R\)，求三棱锥 \(P - ABC\) 的体积.
\end{enumerate}

\bitmapfigure[max width=0.34\linewidth]{img/2008/guangdong_liberal/q18_fig1.png}
\end{problem}

\begin{solution}
\begin{enumerate}
\item 因为 \(BD\) 是圆的直径，所以 \(\angle BAD=90^\circ\). 在直角三角形 \(ABD\) 中，\(BD=2R\)，\(\angle ABD=60^\circ\)，从而 \(AB=R\)，\(AD=\sqrt3R\). 由 \(\triangle ADP\sim\triangle BAD\)，对应关系为 \(A\leftrightarrow B\)，\(D\leftrightarrow A\)，\(P\leftrightarrow D\)，故
\(\frac{AD}{BA}=\frac{DP}{AD}=\sqrt3\). 因此 \(PD=\sqrt3\,AD=3R\).

\item 在直角三角形 \(BCD\) 中，\(BD=2R\)，且 \(\angle BDC=45^\circ\)，所以 \(BC=CD=\sqrt2R\). 由四边形顶点的顺序，\(BA\) 与 \(BC\) 分居直径 \(BD\) 两侧，故 \(\angle ABC=\angle ABD+\angle DBC=60^\circ+45^\circ=105^\circ\). 于是
\(S_{\triangle ABC}=\frac12AB\cdot BC\sin\angle ABC=\frac12\cdot R\cdot\sqrt2R\cdot\sin105^\circ=\frac{\sqrt3+1}{4}R^2\).

又因为 \(PC=\sqrt{11}R\)，且 \(PD=3R\)，\(DC=\sqrt2R\)，有 \(PC^2=PD^2+DC^2\)，所以 \(\angle PDC=90^\circ\). 相似三角形给出 \(\angle ADP=\angle BAD=90^\circ\)，故 \(PD\) 同时垂直于底面内相交的直线 \(AD\)，\(DC\)，从而 \(PD\perp\) 底面 \(ABCD\). 因此三棱锥 \(P-ABC\) 的高为 \(3R\)，体积为
\(V_{P-ABC}=\frac13S_{\triangle ABC}\cdot3R=\frac{\sqrt3+1}{4}R^3\).
\end{enumerate}
\end{solution}

\begin{problem}
某初级中学共有学生 \(2000\) 名，各年级男、女生人数如下表：

\begin{center}
\begin{tabular}{c|ccc}
\hline
 & 一年级 & 二年级 & 三年级 \\
\hline
女生 & \(373\) & \(x\) & \(y\) \\
男生 & \(377\) & \(370\) & \(z\) \\
\hline
\end{tabular}
\end{center}

已知在全校学生中随机抽取 \(1\) 名，抽到初二年级女生的概率是 \(0.19\).

\begin{enumerate}
\item 求 \(x\) 的值；
\item 现用分层抽样的方法在全校抽取 \(48\) 名学生，问应在初三年级抽取多少名？
\item 已知 \(y \geqslant 245\)，\(z \geqslant 245\)，求初三年级中女生比男生多的概率.
\end{enumerate}
\end{problem}

\begin{solution}
\begin{enumerate}
\item 由抽到初二年级女生的概率为 \(0.19\)，得 \(\frac{x}{2000}=0.19\)，所以 \(x=380\).

\item 由全校学生总数，得 \(y+z=2000-(373+377+380+370)=500\). 初三年级学生占全校学生的比例为 \(\frac{500}{2000}\)，所以分层抽取 \(48\) 名学生时，初三年级应抽取 \(48\times\frac{500}{2000}=12\) 名.

\item 因为 \(y+z=500\)，且 \(y\)，\(z\) 为正整数并满足 \(y\geqslant245\)，\(z\geqslant245\)，所以 \(y\) 可以取 \(245\)，\(246\)，\(\ldots\)，\(255\)，共有 \(11\) 个等可能的有序数对 \((y,z)\). 初三年级女生比男生多等价于 \(y>z\)，结合 \(y+z=500\) 即 \(y>250\)，所以有 \(y=251\)，\(252\)，\(253\)，\(254\)，\(255\) 共 \(5\) 个有序数对满足条件. 因而所求概率为 \(\frac5{11}\).
\end{enumerate}
\end{solution}

\begin{problem}
设 \(b > 0\)，椭圆方程为 \(\frac{x^2}{2b^2} + \frac{y^2}{b^2} = 1\)，抛物线方程为 \(x^2 = 8(y - b)\). 如图所示，过点 \(F(0, b + 2)\) 作 \(x\) 轴的平行线，与抛物线在第一象限的交点为 \(G\). 已知抛物线在点 \(G\) 的切线经过椭圆的右焦点 \(F_1\).

\begin{enumerate}
\item 求满足条件的椭圆方程和抛物线方程；
\item 设 \(A\)，\(B\) 分别是椭圆长轴的左，右端点. 试探究在抛物线上是否存在点 \(P\)，使得 \(\triangle ABP\) 为直角三角形；若存在，请指出共有几个这样的点，并说明理由（不必具体求出这些点的坐标）.
\end{enumerate}

\bitmapfigure[max width=0.44\linewidth]{img/2008/guangdong_liberal/q20_fig1.png}
\end{problem}

\begin{solution}
\begin{enumerate}
\item 抛物线可写为 \(x^2=8(y-b)\). 直线 \(y=b+2\) 与抛物线在第一象限的交点满足 \(x^2=16\)，故 \(G(4,b+2)\). 对抛物线方程求导，得 \(2x=8y'\)，所以在 \(G\) 点的切线斜率为 \(1\)，切线方程为 \(y-(b+2)=x-4\)，即 \(y=x+b-2\).

椭圆的长半轴平方为 \(2b^2\)，短半轴平方为 \(b^2\)，所以焦距的一半满足 \(c^2=2b^2-b^2=b^2\). 因 \(b>0\)，椭圆右焦点为 \(F_1(b,0)\). 由切线经过 \(F_1\)，得 \(0=b+b-2\)，即 \(b=1\). 因此椭圆方程为 \(\frac{x^2}{2}+y^2=1\)，抛物线方程为 \(x^2=8(y-1)\).

\item 此时椭圆长轴的左、右端点分别为 \(A(-\sqrt2,0)\)，\(B(\sqrt2,0)\)，抛物线上点 \(P\) 可表示为 \(P\left(x,1+\frac{x^2}{8}\right)\).

若 \(\angle PAB=90^\circ\)，由于 \(AB\) 水平，必有 \(x=-\sqrt2\)，抛物线上恰有一个这样的点. 同理，若 \(\angle PBA=90^\circ\)，必有 \(x=\sqrt2\)，也恰有一个这样的点.

若 \(\angle APB=90^\circ\)，由 \(\overrightarrow{PA}\cdot\overrightarrow{PB}=0\)，得
\(x^2-2+\left(1+\frac{x^2}{8}\right)^2=0\). 令 \(t=x^2\)，则 \(t^2+80t-64=0\)，解得 \(t=-40\pm8\sqrt{26}\). 其中只有 \(t=-40+8\sqrt{26}>0\) 符合，因此 \(x\) 有两个相反的取值，对应两个点. 综上，抛物线上共有 \(1+1+2=4\) 个点，使 \(\triangle ABP\) 为直角三角形.
\end{enumerate}
\end{solution}

\begin{problem}
设数列 \(\{a_n\}\) 满足 \(a_1 = 1\)，\(a_2 = 2\)，\(a_n = \frac{1}{3}(a_{n-1} + 2a_{n-2})\)（\(n = 3\)，\(4\)，\(\cdots\)）. 数列 \(\{b_n\}\) 满足 \(b_1 = 1\)，\(b_n\)（\(n = 2\)，\(3\)，\(\cdots\)）是非零整数，且对任意的正整数 \(m\) 和自然数 \(k\)，都有 \(-1 \leqslant b_m + b_{m+1} + \cdots + b_{m+k} \leqslant 1\).

\begin{enumerate}
\item 求数列 \(\{a_n\}\) 和 \(\{b_n\}\) 的通项公式；
\item 若 \(c_n = na_nb_n\)（\(n = 1\)，\(2\)，\(\dots\)），求数列 \(\{c_n\}\) 的前 \(n\) 项和 \(S_n\).
\end{enumerate}
\end{problem}

\begin{solution}
\begin{enumerate}
\item 数列 \(\{a_n\}\) 的特征方程为 \(3r^2-r-2=0\)，即 \((r-1)(3r+2)=0\)，所以特征根为 \(1\) 和 \(-\frac23\). 设 \(a_n=A+B\left(-\frac23\right)^{n-1}\)，由 \(a_1=1\)，\(a_2=2\) 得
\(A+B=1\)，\(A-\frac23B=2\). 解得 \(A=\frac85\)，\(B=-\frac35\)，故
\(a_n=\frac85-\frac35\left(-\frac23\right)^{n-1}\).

取题设中的 \(k=0\)，得 \(-1\leqslant b_m\leqslant1\). 由于 \(b_m\) 为非零整数，所以 \(b_m\in\{-1,1\}\). 再取 \(k=1\)，相邻两项的和不能等于 \(2\) 或 \(-2\)，故 \(b_{m+1}=-b_m\). 结合 \(b_1=1\)，得到 \(b_n=(-1)^{n+1}\).

\item 由上式，
\(c_n=na_nb_n=\frac85n(-1)^{n+1}-\frac35n\left(\frac23\right)^{n-1}\). 因而
\(S_n=\frac85\sum_{k=1}^n k(-1)^{k+1}-\frac35\sum_{k=1}^n k\left(\frac23\right)^{k-1}\).

由等比数列求和公式及其求导形式，
\(\displaystyle\sum_{k=1}^n k(-1)^{k+1}=\frac{1+(-1)^{n+1}(2n+1)}4\)，以及
\(\displaystyle\sum_{k=1}^n kr^{k-1}=\frac{1-(n+1)r^n+nr^{n+1}}{(1-r)^2}\). 令 \(r=\frac23\)，代入并化简，得
\(S_n=-5+\frac25(-1)^{n+1}(2n+1)+\frac95(n+3)\left(\frac23\right)^n\).
\end{enumerate}
\end{solution}
